% Auto-generated by 1_code/7_main_analysis/2_appendix/b2_semantic_gap_text_interpretability.py — do not edit manually
\begin{tabular}{lcp{0.26\textwidth}p{0.26\textwidth}p{0.25\textwidth}}
\toprule
SDG & Gap & Research-side distinctive terms & Policy-side distinctive terms & Descriptive reading \\
\midrule
17 & 0.219 & sentiment, agricultural, tweets, used, fusion, covid-, agriculture, twitter & countries, international, nations, oecd, united, agenda, global, report & High-gap case: the research-side sample is technical and biomedical/ML-heavy, while policy language is international, institutional, and agenda-oriented. \\
13 & 0.478 & soil, crop, agricultural, yield, moisture, rice, used, stress & climate, change, emissions, national, countries, action, global, energy & High-gap case: both sides concern climate, but policy language is more institutional and commitment-oriented while research language is more technical and measurement-oriented. \\
9 & 0.289 & agricultural, agriculture, iot, cancer, food, crop, detection, cell & infrastructure, innovation, national, government, public, population, countries, access & Lower-gap comparison: differences remain visible, but the contrast is closer to a technology/infrastructure register split than to the broader institutional gap seen in SDGs 13 and 17. \\
\bottomrule
\end{tabular}
